\chapter{Differentialregning}
\label{chap:ch10}

Både i praksis og i moderne ingeniørarbejde møder vi ofte situationer, hvor én størrelse ændrer sig som følge af en anden. En bils brændstofforbrug afspejler, hvordan den tilbagelagte afstand afhænger af den mængde brændstof, der bruges. En sprinters hastighed beskriver, hvordan positionen varierer med tiden. I landbruget ændrer udbyttet sig med mængden af gødning, og i økonomi reagerer efterspørgslen på prisen. Studiet af sådanne sammenhænge er kernen i \emph{differentialregning}, den matematiske disciplin, der handler om at forstå, hvordan størrelser varierer i forhold til hinanden.

Det begreb, der indfanger denne variation, er \emph{ændringshastigheden}.

\begin{itemize}
\item Ændringen i afstand i forhold til tid er \emph{hastighed}.
\item Ændringen i fortjeneste i forhold til pris er \emph{marginal fortjeneste}.
\item Ændringen i tab i forhold til modelparametre er \emph{gradienten}.
\end{itemize}

Når en ændringshastighed er konstant, er sammenhængen mellem størrelserne lineær, og grafen er en ret linje. Når ændringshastigheden varierer, bliver grafen krum, og hældningen afhænger af det konkrete punkt, hvor vi måler. Det fører naturligt til idéen om en \emph{tangent}, som lokalt approksimerer en kurve i et enkelt punkt.

% --- Ensure these packages are in your preamble ---
% \usepackage{pgfplots}
% \usepackage{xcolor}
% \pgfplotsset{compat=1.17}

% --- Ensure your custom color is defined ---
% \definecolor{mseViaBlue}{HTML}{1933CC}

\begin{figure}[htbp]
    \centering
    \begin{tikzpicture}
        \begin{axis}[
            axis lines=middle,
            xlabel=$x$,
            ylabel=$y$,
            width=12cm, height=8cm,
            xmin=-0.5, xmax=2.5,
            ymin=-1, ymax=9,
            xtick={-1, 0, 1, 2},
            ytick={0, 2, 4, 6, 8},
            grid=major,
            grid style={dashed, gray!30},
            legend style={
                at={(0.5, 0.97)}, % Position near the top center
                anchor=north,     % Anchor the top-center of the legend box
                legend cell align={left}
            },
            clip=false
        ]
        % Plot the curve f(x) = x^3
        \addplot[
            domain=0:2.1, samples=100,
            smooth, very thick, mseViaBlue
        ] {x^3};
        \addlegendentry{\(f(x) = x^3\)}

        % Plot the tangent line at x=1, which is y = 3x - 2
        \addplot[
            domain=0:2.5,
            red, thick
        ] {3*x - 2};
        \addlegendentry{Tangent i \(x=1\)}

        % Mark the point of tangency
        \coordinate (P) at (axis cs:1, 1);
        \fill[red] (P) circle (2.5pt);
        \node[below right, font=\small] at (P) {\(P(1, 1)\)};
        
        \end{axis}
    \end{tikzpicture}
    \caption{Tangenten til grafen for \(f(x)=x^3\) i punktet \(P(1, 1)\).}
    \label{fig:ch10-tangent}
\end{figure}

\begin{remark}
For en ret linje er hældningen konstant. For en krum graf varierer hældningen fra punkt til punkt, og vi taler derfor om den \emph{momentane ændringshastighed}, dvs. tangentens hældning i et bestemt punkt.
\end{remark}

Det samme princip gælder for beregningsmæssige og datadrevne systemer. Et programs køretid ændrer sig med inputstørrelsen, en models nøjagtighed varierer med antallet af træningsepoker, og et netværks gennemløb afhænger af den tilgængelige båndbredde. Inden for maskinlæring minder træningen af en model om en vandrer, der leder efter det laveste punkt i et stort, tågedækket landskab. »Højden« repræsenterer modellens fejl, og for at komme videre må man gå nedad. Hvilken retning der er den rigtige blandt millioner af mulige, afhænger af, at vi kender den lokale hældning -- den \emph{afledte} -- som udtrykker den momentane ændringshastighed.

Differentialregning giver os værktøjerne til præcist at beskrive, beregne og fortolke sådanne ændringer. Den omsætter intuitive forestillinger som at være hurtigere, stejlere eller mere effektiv til eksakte matematiske udtryk. Dette kapitel introducerer de grundlæggende idéer i differentialregning. Vi begynder med begreberne \emph{grænseværdi} og \emph{kontinuitet}, som formaliserer forestillingen om at nærme sig et punkt uendeligt tæt. På det grundlag udvikler vi definitionen og fortolkningen af den \emph{afledte} som et mål for momentan ændring og ser på anvendelser i sammenhænge, der spænder fra fysisk bevægelse til optimering af moderne algoritmer.


\section{Grænseværdier og kontinuitet: grundlaget for differentialregning}

Før vi kan måle en momentan ændringshastighed, må vi først opbygge en formel forståelse af, hvad det vil sige at komme »uendeligt tæt« på et punkt uden nødvendigvis at være i punktet. Det er begrebet \textbf{grænseværdi}. Med udgangspunkt i grænseværdier giver \textbf{kontinuitet} os et sprog til at beskrive funktioner, der er forudsigelige og velopdragne, uden pludselige spring eller brud.

\subsection*{Grænseværdier}

En grænseværdi beskriver den værdi, som en funktion »nærmer sig«, når dens input kommer tættere og tættere på et bestemt punkt. Idéen er grundlæggende ikke kun i differentialregning, men også i softwareudvikling, hvor vi ofte analyserer systemers opførsel, når forholdene nærmer sig en bestemt tilstand.

Når vi for eksempel analyserer en algoritmes køretid, spørger vi, hvad der sker med udførelsestiden, når inputstørrelsen `n` går mod uendelig. Inden for netværk kan vi modellere det teoretiske maksimale gennemløb som en grænseværdi, når latensen går mod nul.

\begin{definition}[Grænseværdi]
Lad \( f(x) \) være en funktion, der er defineret i nærheden af et punkt \( c \). Vi siger, at \textbf{grænseværdien} af \( f(x) \), når \( x \) går mod \( c \), er \( L \), og skriver
\vspace{0.5em}
\[
\lim_{x \to c} f(x) = L
\]
\vspace{0.5em}
hvis vi kan gøre værdien af \( f(x) \) vilkårligt tæt på \( L \) ved at vælge et \( x \), der er tilstrækkeligt tæt på \( c \), men forskelligt fra \( c \).
\end{definition}

I mange simple tilfælde er grænseværdien blot funktionens værdi i punktet. For \( f(x) = x^2 \) er grænseværdien, når \( x \) går mod 3, for eksempel simpelthen \( 3^2 = 9 \). Grænseværdiernes egentlige styrke viser sig dog i tilfælde, hvor funktionen ikke er defineret i det punkt, vi interesserer os for.

\begin{example} En grænseværdi i et punkt, hvor funktionen ikke er defineret

Se på funktionen \( f(x) = \frac{\sin(x)}{x} \). Funktionen er grundlæggende i signalbehandling og kaldes sinc-funktionen. Bemærk, at \( f(0) \) ikke er defineret, fordi det giver division med nul. Vi kan dog spørge, hvilken værdi funktionen nærmer sig, når \( x \) kommer meget tæt på 0.

Ved at indsætte værdier meget tæt på 0 kan vi se en tendens:
\begin{align*}
f(0.1) &= \frac{\sin(0.1)}{0.1} \approx 0.9983 \\
f(0.01) &= \frac{\sin(0.01)}{0.01} \approx 0.99998 \\
f(-0.01) &= \frac{\sin(-0.01)}{-0.01} \approx 0.99998
\end{align*}
Funktionen ser ud til at nærme sig 1. I differentialregning kan det bevises formelt, at:
\[
\lim_{x \to 0} \frac{\sin(x)}{x} = 1
\]
Evnen til at analysere opførslen i et problematisk punkt er afgørende for at kunne lave robuste numeriske algoritmer.
\end{example}

\begin{example} Asymptotisk analyse med grænseværdier

Antag, at du vil undersøge, om funktionen \( f(n) = 5n^2 + 4n + 7 \) er \( \mathcal{O}(n^2) \), når \( n \) bliver stor, sådan som man ofte gør i studiet af algoritmers kompleksitet. En måde at formalisere det på er at undersøge grænseværdien:
\[
\lim_{n \to \infty} \frac{f(n)}{n^2}
= \lim_{n \to \infty} \frac{5n^2 + 4n + 7}{n^2}
= \lim_{n \to \infty} \left( 5 + \frac{4}{n} + \frac{7}{n^2} \right)
\]
Når \( n \) bliver meget stor, går leddene \( \frac{4}{n} \) og \( \frac{7}{n^2} \) begge mod nul, så hele udtrykket nærmer sig \( 5 \). Da grænseværdien er en endelig konstant, bekræfter det, at \( f(n) \) vokser med samme hastighed som \( n^2 \). Denne type grænseværdiberegning er kernen i asymptotisk analyse og giver en formel begrundelse for påstande om algoritmers køretid udtrykt med store \(\mathcal{O}\).
\end{example}

\subsection*{Kontinuitet}
Kontinuitet formaliserer idéen om, at en funktion er »velopdragen« eller forudsigelig. Intuitivt er en funktion kontinuert, hvis du kan tegne dens graf uden at løfte blyanten fra papiret. I software antager vi ofte, at vores systemer er kontinuerte: En lille ændring i input bør give en lille, forudsigelig ændring i output og ikke et pludseligt, drastisk spring.

\begin{definition}[Kontinuitet]
En funktion \( f \) er \textbf{kontinuert} i et punkt \( c \), hvis følgende tre betingelser er opfyldt:
\begin{enumerate}
    \item \( f(c) \) er defineret.
    \item \( \lim_{x \to c} f(x) \) eksisterer.
    \item \( \lim_{x \to c} f(x) = f(c) \).
\end{enumerate}
En funktion er kontinuert på et interval, hvis den er kontinuert i alle punkter i intervallet.
\end{definition}

Diskontinuiteter repræsenterer ofte vigtige hændelser i softwaresystemer, fx tilstandsskift eller tærskler, der udløser en handling.

\begin{example} Springdiskontinuitet i en SaaS-prismodel
    
Se på en cloudtjeneste, der tager betaling efter antallet af API-kald pr. måned. Prismodellen er trinvis: \$20 for op til 10\,000 kald og \$40 for ethvert forbrug over 10\,000 kald. Omkostningsfunktionen \(C(x)\), hvor \(x\) er antallet af kald, har en \textbf{springdiskontinuitet} i \(x = 10\,000\).
\[
C(x) =
\begin{cases}
    20 & \text{hvis } x \le 10000 \\
    40 & \text{hvis } x > 10000
\end{cases}
\]
Som vist i \autoref{fig:jump_discontinuity} springer prisen pludselig ved tærsklen. Det er afgørende at forstå diskontinuiteter, når man modellerer systemer med tærskelbaseret logik, lige fra regler for automatisk skalering til faktureringssystemer.
\end{example}

\begin{figure}[htbp]
\centering
\begin{tikzpicture}
\begin{axis}[
    xlabel={API-kald},
    ylabel={Pris (\$)},
    xmin=0, xmax=20,
    ymin=0, ymax=50,
    xtick={0, 10, 20},
    xticklabels={0, 10k, 20k},
    ytick={0, 20, 40},
    grid=major,
    grid style={dashed,gray!30},
    title style={font=\bfseries}
]
% First tier: C(x) = 20 for x <= 10000
\addplot[mseViaBlue, very thick, domain=0:10] {20};
% Second tier: C(x) = 40 for x > 10000
\addplot[mseViaBlue, very thick, domain=10:20] {40};

% Mark the discontinuity
% Solid circle at the end of the first tier
\addplot[only marks, mark=*, mark size=2pt, mseViaBlue] coordinates {(10, 20)};
% Hollow circle at the start of the second tier
\addplot[only marks, mark=o, mark size=2pt, mseViaBlue, thick] coordinates {(10, 40)};
\end{axis}
\end{tikzpicture}
\caption{Omkostningsfunktionen har en springdiskontinuitet ved 10\,000 API-kald, hvor pristrinnet skifter.}
\label{fig:jump_discontinuity}
\end{figure}

\section{Den afledte: et mål for ændringshastigheden}
\label{sec:ch10-derivative}

Vi taler ofte om gennemsnitlige ændringshastigheder. Hvis en dataoverførsel tager 10 sekunder og flytter 500 MB, er den gennemsnitlige overførselshastighed for eksempel 50 MB/s. Men gennemsnittet fortæller intet om udsvingene under overførslen. Den \textbf{afledte} er det værktøj, der lader os gå fra den gennemsnitlige hastighed til den \textbf{momentane ændringshastighed} på et hvilket som helst bestemt tidspunkt.

Vi kan altså sige, at differentialregning handler om at bestemme, hvordan én størrelse ændrer sig i forhold til en anden. Geometrisk er den afledte af en funktion i et punkt hældningen af tangenten til funktionens graf i punktet. For at finde en tangents hældning approksimerer vi den med sekanter gennem to nærliggende punkter på kurven. Når punkterne rykker tættere på hinanden, nærmer sekantens hældning sig tangentens hældning. Som vist i \autoref{fig:secant_to_tangent} er tangenten grænsen for sekanter gennem to punkter på kurven, når punkterne kommer uendeligt tæt på hinanden.

\begin{figure}[htbp]
    \centering
    \begin{tikzpicture}
    \begin{axis}[
        axis lines=middle, xlabel={$x$}, ylabel={$y$},
        xtick=\empty, ytick=\empty,
        xmin=-0.5, xmax=3, ymin=-0.5, ymax=9,
        width=14cm, height=10cm,
        legend style={at={(0.85,1)}, anchor=north west, draw=none, fill=none, legend cell align=left},
        title style={font=\bfseries}
    ]
    \addplot[domain=0:2.5, samples=150, smooth, very thick, mseViaBlue] {x^3};
    \addlegendentry{\(f(x) = x^3\)}
    \coordinate (P) at (axis cs:1, 1);
    \coordinate (Q2) at (axis cs:1.5, 3.375);
    \coordinate (Q1) at (axis cs:2, 8);
    \addplot[domain=0.6:2.2, dashed, color=gray] {7*x - 6};
    \addlegendentry{Sekant gennem $PQ_1$}
    \addplot[domain=0.6:2.0, color=orange, dashed] {4.75*x - 3.75};
    \addlegendentry{Sekant gennem $PQ_2$}
    \addplot[domain=0.5:2.0, red] {3*x - 2};
    \addlegendentry{Tangent i $P$}
    \fill[red] (P) circle (2pt) node[below right, xshift=4pt] {$P$};
    \fill[black] (Q2) circle (2pt) node[right, xshift=4pt] {$Q_2$};
    \fill[black] (Q1) circle (2pt) node[above right, xshift=4pt] {$Q_1$};
    \end{axis}
    \end{tikzpicture}
    \caption{Tangenten som grænse for sekanter.}
    \label{fig:secant_to_tangent}
    \end{figure}
    
    Når punktet $Q$ i \autoref{fig:secant_to_tangent} bevæger sig langs kurven mod $P$ (fra $Q_1$ til $Q_2$ til $P$), nærmer sekanten $PQ$'s hældning sig hældningen af tangenten til $f(x) = x^3$ i $P$. Bemærk, at tangenten kun rører kurven i $P$, mens sekanten går gennem to punkter.
    
    Det fører til den formelle definition af den afledte.
    
    \begin{definition}[Den afledte]
    Den \textbf{afledte} af en funktion \( f \) med hensyn til \( x \), betegnet \( f'(x) \), er funktionen
    \[
    f'(x) = \lim_{h \to 0} \frac{f(x+h) - f(x)}{h}
    \]
    forudsat at grænseværdien eksisterer. Hvis \( f'(x) \) eksisterer, siger vi, at \( f \) er \textbf{differentiabel} i \( x \).
    \end{definition}
    
    Udtrykket \( \frac{f(x+h) - f(x)}{h} \) er \textbf{differenskvotienten}, som er den gennemsnitlige ændringshastighed på intervallet fra \(x\) til \(x+h\). Den afledte er grænseværdien af denne gennemsnitlige hastighed, når intervallets længde \(h\) går mod nul.
    
    \begin{example} Den afledte ud fra definitionen
        
    Lad os finde den afledte af \( f(x) = x^2 \) ved hjælp af definitionen med grænseværdi.
    \end{example}
    \begin{solution}
    Vi begynder med differenskvotienten og reducerer:
    \[
    \frac{f(x+h) - f(x)}{h} = \frac{(x+h)^2 - x^2}{h} = \frac{x^2 + 2xh + h^2 - x^2}{h} = \frac{2xh + h^2}{h} = 2x + h
    \]
    Til sidst tager vi grænseværdien for \( h \to 0 \):
    \[
    f'(x) = \lim_{h \to 0} (2x + h) = 2x
    \]
    Den afledte af \( f(x) = x^2 \) er altså \( f'(x) = 2x \). Den angiver parablens hældning i ethvert punkt \(x\).
    \end{solution}

\subsection*{Notation}
Den afledte af $f$ kan skrives på flere ækvivalente måder:
\[
f'(x), \quad \frac{df}{dx}, \quad \frac{dy}{dx}, \quad Df(x), \quad \frac{d}{dx}f(x).
\]

\section{Differentiationsregler}
\label{sec:ch10-rules}

Differentialregning giver os en metode til at finde den eksakte afledte af en funktion direkte ud fra dens forskrift uden at skulle bruge grafer eller numeriske skøn. I praksis sker det ved hjælp af en række enkle regler, som gør det muligt at beregne den afledte af næsten enhver funktion, vi sandsynligvis vil møde. I dette afsnit introducerer vi reglerne, forklarer deres betydning og viser, hvordan de anvendes i praksis.


\subsection*{Konstantreglen og potensreglen}
De måske simpleste funktioner i matematikken er de konstante funktioner og funktionerne på formen $x^{n}$.\\

\begin{theorem}{Konstantreglen og potensreglen}
    \begin{enumerate}
        \item \textbf{Konstantreglen:} Hvis $c$ er en konstant, så er $\dfrac{d}{dx}(c) = 0$, \vspace{0.5em}
        \item \textbf{Potensreglen:} For ethvert reelt $n$ gælder $\dfrac{d}{dx}x^n = n x^{n - 1}$. \vspace{0.5em}
    \end{enumerate}
\end{theorem}

\begin{example}

    \begin{itemize}
        \item Hvis \(f(x)=x^7\), så er \(f'(x) = 7x^6\),
        \vspace{0.5em}
        \item Hvis \(y = x^{-0.5}\), så er \(\frac{dy}{dx} = -0.5x^{-1.5}\),
        \vspace{0.5em}
        \item \(\frac{d}{dx} x^{-3} = -3x^{-4}\),
        \vspace{0.5em}
        \item Hvis \(g(x) = 3.2\), så er \(g'(x) = 0\),
        \vspace{0.5em}
        \item Hvis \(f(t) = \sqrt{t}=t^{\frac{1}{2}}\), så er \(f'(t) = \frac{1}{2} t^{-\frac{1}{2}}=\frac{1}{2 \sqrt{t}} \),
        \vspace{0.5em}
        \item Hvis \(h(u) = -13.29\), så er \(h'(u) = 0\),
        \vspace{0.5em}
        \item $\dfrac{d}{dx}x^3 = 3x^2$,
        $\dfrac{d}{dx}x^{-2} = -2x^{-3}$.
    \end{itemize}

\end{example}

I eksemplerne ovenfor har vi brugt konstantreglen og potensreglen til at finde de afledte af flere simple funktioner. Det er dog vigtigt at huske, hvad disse beregninger faktisk repræsenterer. Den afledte angiver hældningen af tangenten til en funktions graf i et givet punkt.

Hvis for eksempel $f(x) = x^{2}$, så er $f'(x) = 2x$. For at finde hældningen af tangenten til grafen for $x^{2}$ i $x = 1$ indsætter vi $x = 1$ i den afledte og får $f'(1) = 2 \times 1 = 2$. Tilsvarende er tangentens hældning i $x = -0.5$ lig med $f'(-0.5) = 2 \times (-0.5) = -1$.

Resultaterne er illustreret i \autoref{fig:tangent_lines}, hvor tangenterne i $x = 1$ og $x = -0.5$ har hældningerne henholdsvis 2 og -1.

% --- Ensure these packages are in your preamble ---
% \usepackage{pgfplots}
% \usepackage{xcolor}
% \pgfplotsset{compat=1.17}

% --- Ensure your custom color is defined ---
% \definecolor{mseViaBlue}{HTML}{1933CC}

\begin{figure}[htbp]
    \centering
    \begin{tikzpicture}
        \begin{axis}[
            axis lines=middle,
            xlabel=$x$,
            ylabel=$y$,
            width=12cm, height=8cm,
            xmin=-1.5, xmax=1.5,
            ymin=-0.5, ymax=2.5,
            xtick={-1.5, -1, -0.5, 0, 0.5, 1, 1.5},
            ytick={0, 0.5, 1, 1.5, 2, 2.5},
            grid=major,
            grid style={dashed, gray!30},
            legend pos=outer north east,
            legend cell align={left},
            clip=false,
        ]
        % Plot the curve f(x) = x^2
        \addplot[
            domain=-1.5:1.5, samples=100,
            smooth, very thick, mseViaBlue
        ] {x^2};
        \addlegendentry{\(f(x)=x^2\)}

        % -- Tangent Line at x = 1 --
        \addplot[
            domain=0:1.5,
            red, thick
        ] {2*x - 1};
        \addlegendentry{Tangent i \(x=1\)}
        
        % Mark and label the point at x=1
        \coordinate (P1) at (axis cs:1, 1);
        \fill[red] (P1) circle (2.5pt);
        \node[right, xshift=10pt, font=\small, color=red] at (P1) {Hældning \(f'(1)=2\)};

        % -- Tangent Line at x = -0.5 --
        \addplot[
            domain=-1.5:0.5,
            orange, thick
        ] {-x - 0.25};
        \addlegendentry{Tangent i \(x=-0.5\)}

        % Mark and label the point at x=-0.5
        \coordinate (P2) at (axis cs:-0.5, 0.25);
        \fill[orange] (P2) circle (2.5pt);
        \node[left, xshift=-10pt, font=\small, color=orange] at (P2) {Hældning \(f'(-0.5)=-1\)};

        \end{axis}
    \end{tikzpicture}
    \caption{Tangenterne til grafen for \(f(x) = x^2\) har forskellige hældninger i forskellige punkter.}
    \label{fig:tangent_lines}
\end{figure}


\begin{example}

Lad $f(x)=(x^2+1)^3$. Så er
\[
f'(x) = 3(x^2+1)^2\cdot 2x = 6x(x^2+1)^2.
\]
\end{example}

\begin{example}

Find hældningen af tangenten til grafen for funktionen $g(t)=t^{4}$ i det punkt på grafen, hvor $t=-2$.


\begin{solution}
    
Den afledte er $g^{\prime}(t)=4 t^{3}$, så tangentens hældning i $t=-2$ er $g^{\prime}(-2)= 4 \times(-2)^{3}=-32$.
\end{solution}

\end{example}



\begin{example}

Find ligningen for tangenten til grafen for $y=f(x)=x^{\frac{1}{2}}$ i punktet $x=4$.



\begin{solution}

$f(4)=4^{\frac{1}{2}}=\sqrt{4}=2$, så koordinaterne til punktet på grafen er $(4,2)$. Den afledte er

Enhver ikke-lodret linje har en ligning på formen $y=m x+b$, hvor $m$ er hældningen, og $b$ er skæringen med andenaksen.

Her er hældningen $\frac{1}{4}$, så $m=\frac{1}{4}$, og ligningen er $y=\frac{x}{4}+b$. Da linjen går gennem punktet $(4,2)$, ved vi, at $y=2$, når $x=4$.

Ved indsættelse får vi $2=\frac{4}{4}+b$, så $b=1$. Ligningen er derfor $y=\frac{x}{4}+1$. \end{solution}

\end{example}

Vi ved nu, hvordan man differentierer enhver funktion, der er en potens af variablen. Eksempler er funktioner som $x^{3}$ og $t^{-1.3}$. Du vil møde funktioner, der ved første øjekast ikke ser ud til at være en potens af variablen, men som kan omskrives til denne form. Et af de simpleste eksempler er funktionen

\[
    f(t)=\sqrt{t}
\]

som også kan skrives på formen

\[
f(t)=t^{\frac{1}{2}} .
\]

Den afledte er så

\[
f^{\prime}(t)=\frac{t^{-\frac{1}{2}}}{2}=\frac{1}{2 \sqrt{t}}
\]

Tilsvarende gælder, at hvis

\[
h(s)=\frac{1}{s}=s^{-1}
\]

så er

\[
h^{\prime}(s)=-s^{-2}=-\frac{1}{s^{2}}
\]

\begin{example}

Hvis $f(x)=\frac{1}{\sqrt[3]{x}}=x^{-\frac{1}{3}}$, så er $f^{\prime}(x)=-\frac{1}{3} x^{-\frac{4}{3}}$.
\end{example}

\begin{example}

Hvis $y=\frac{1}{x \sqrt{x}}=x^{-\frac{3}{2}}$, så er $\frac{d y}{d x}=-\frac{3}{2} x^{-\frac{5}{2}}$.
\end{example}


\subsection*{Skalarreglen og linearitet}

Indtil videre ved vi, hvordan man differentierer potenser af den uafhængige variabel. Mange af de funktioner, der optræder i anvendelser, er bygget op af sådanne potenser på enkle måder. Funktionen $3x^2$ er for eksempel blot et skalart multiplum af $x^2$, men hverken konstantreglen eller potensreglen fortæller os, hvordan vi differentierer $3x^2$. Reglerne forklarer heller ikke, hvordan man differentierer udtryk som $x^{2} + x^{3}$ eller $x^{2} - x^{3}$.

\begin{theorem}{Skalarreglen og linearitetsreglen}
    \begin{enumerate}
        \setcounter{enumi}{2}
        \vspace{0.5em}
        \item \textbf{Skalarreglen:} \( \frac{d}{dx}[c \cdot f(x)] = c \cdot f'(x) \). \vspace{0.5em}
        \item \textbf{Linearitetsreglen:} \( \frac{d}{dx}[f(x) \pm g(x)] = f'(x) \pm g'(x) \). \vspace{0.5em}
    \end{enumerate}
\end{theorem}

Skalarreglen og linearitetsreglen udfylder dette hul. De beskriver, hvordan man differentierer funktioner, der fremkommer ved at gange en funktion med en konstant og ved at lægge funktioner sammen eller trække dem fra hinanden. Med disse regler kan vi udvide vores differentiationsteknikker fra simple potenser til en bred klasse af funktioner, der er bygget op af dem.

\begin{example}
    \begin{itemize}
        \item Hvis \(f(x)=3 x^2\), så er \(f^{\prime}(x)=3 \times \frac{d}{d x} x^2=6 x\). \vspace{0.5em}
        \item Hvis \(g(t)=3 t^2+2 t^{-2}\), så er \(g^{\prime}(t)=\frac{d}{d t} 3 t^2+\frac{d}{d t} 2 t^{-2}=6 t-4 t^{-3}\). \vspace{0.5em}
        \item Hvis \(y=\frac{3}{\sqrt{x}}-2 x \sqrt[3]{x}=3 x^{-\frac{1}{2}}-2 x^{\frac{4}{3}}\), så er \(\frac{d y}{d x}=-\frac{3}{2} x^{-\frac{3}{2}}-\frac{8}{3} x^{\frac{1}{3}}\). \vspace{0.5em}
        \item Hvis \(y=-0.3 x^{-0.4}\), så er \(\frac{d y}{d x}=0.12 x^{-1.4}\). \vspace{0.5em}
        \item \(\frac{d}{d x} 2 x^{0.3}=0.6 x^{-0.7}\).
    \end{itemize}
    
\end{example}

\textbf{Advarsel.} Selv om linearitetsreglen siger, at $\dfrac{d}{dx}\big(f(x)\pm g(x)\big)=f'(x)\pm g'(x)$, gælder denne egenskab \emph{ikke} for produkter eller kvotienter. Generelt er
\[
\frac{d}{dx}\big(f(x)g(x)\big)\neq f'(x)g'(x),
\qquad
\frac{d}{dx}\left(\frac{f(x)}{g(x)}\right)\neq \frac{f'(x)}{g'(x)}.
\]
For at differentiere $f(x)g(x)$ eller $\frac{f(x)}{g(x)}$ kan vi ikke blot differentiere faktorerne og derefter gange eller dividere resultaterne. De korrekte teknikker er \emph{produktreglen} og \emph{kvotientreglen}, som udvikles i næste afsnit.

\subsection*{Produktreglen og kvotientreglen}

En anden almindelig måde at kombinere funktioner på er at gange dem med hinanden og derved danne et \emph{produkt}. Produktreglen giver metoden til at differentiere funktioner, der er opbygget på denne måde. 

Kvotientreglen gør det derimod muligt at differentiere funktioner, der fremkommer ved at dividere én funktion
med en anden, dvs. ved at danne kvotienter af funktioner. Kvotientreglen giver metoden til at differentiere funktioner, der er opbygget på denne måde.

\begin{theorem}{Produktreglen og kvotientreglen}
    \begin{enumerate}
        \setcounter{enumi}{4}
        \vspace{0.5em}
        \item \textbf{Produktreglen:} Hvis \(f(x)\) og \(g(x)\) er differentiable, så gælder \vspace{0.5em}
        \[
        \frac{d}{dx}\big(f(x)g(x)\big)=f'(x)g(x)+f(x)g'(x). \vspace{0.5em}
        \] \vspace{0.5em}
        \item \textbf{Kvotientreglen:} Hvis \(f(x)\) og \(g(x)\) er differentiable, og \(g(x)\neq0\), så gælder
        \vspace{0.5em}
        \[
        \frac{d}{dx}\left(\frac{f(x)}{g(x)}\right)=\frac{f'(x)g(x)-f(x)g'(x)}{[g(x)]^2}. \vspace{0.5em}
        \]
    \end{enumerate}
\end{theorem}

\begin{example} Produktreglen

    \begin{itemize}
        \item Hvis \(y=(x+2)\left(x^2+3\right)\), så er \(y^{\prime}=(x+2) 2 x+1\left(x^2+3\right)\). \vspace{0.5em}
        \item Hvis \(f(x)=\sqrt{x}\left(x^3-3 x^2+7\right)\), så er \(f^{\prime}(x)=\sqrt{x}\left(3 x^2-6 x\right)+\frac{1}{2} x^{-\frac{1}{2}}\left(x^3-3 x^2+7\right)\). \vspace{0.5em}
        \item Hvis \(z=\left(t^2+3\right)\left(\sqrt{t}+t^3\right)\), så er \(\frac{d z}{d t}=\left(t^2+3\right)\left(\frac{1}{2} t^{-\frac{1}{2}}+3 t^2\right)+2 t\left(\sqrt{t}+t^3\right)\). \vspace{0.5em}
    \end{itemize}

\end{example}

\begin{example} Kvotientreglen

    \begin{itemize}
        \item Hvis \(y=\frac{2 x^2+3 x}{x^3+1}\), så er \(\frac{d y}{d x}=\frac{\left(x^3+1\right)(4 x+3)-\left(2 x^2+3 x\right) 3 x^2}{\left(x^3+1\right)^2}\). \vspace{0.5em}   
        \item Hvis \(g(t)=\frac{t^2+3 t+1}{\sqrt{t}+1}\), så er \(g^{\prime}(t)=\frac{(\sqrt{t}+1)(2 t+3)-\left(t^2+3 t+1\right)\left(\frac{1}{2} t^{-\frac{1}{2}}\right)}{(\sqrt{t}+1)^2}\). \vspace{0.5em}
    \end{itemize}
\end{example}

På grund af minustegnet i tælleren (dvs. øverst i brøken) er det i kvotientreglen vigtigt at skrive leddene i tælleren i den rigtige rækkefølge. Det er en hyppig kilde til fejl, så vær omhyggelig. Find din egen måde at huske den rigtige rækkefølge af leddene på.


\subsection*{Kædereglen}
Kædereglen er nok den vigtigste differentiationsregel, især inden for maskinlæring. Den fortæller os, hvordan vi finder den afledte af en sammensat funktion -- en funktion, der er indlejret i en anden. I software sammensætter vi hele tiden funktioner, og kædereglen er det matematiske værktøj til at analysere, hvordan ændringer forplanter sig gennem disse sammensætninger.

\begin{theorem}{Kædereglen}
    \vspace{0.5em}
    \setcounter{enumi}{6}
    \begin{enumerate}
        \item Hvis \( h(x) = f(g(x)) \) er en sammensat funktion, er dens afledte lig med den afledte af den ydre funktion (udregnet i den indre funktion) ganget med den afledte af den indre funktion. \vspace{0.5em}
        \[
        h'(x) = f'(g(x)) \cdot g'(x). \vspace{0.5em}
        \]
    \end{enumerate}
\end{theorem}

\begin{example}

Find den afledte af \( h(x) = (x^3 + 2x)^5 \).
\end{example}

\begin{solution}
Der er tale om en sammensætning, hvor den ydre funktion er \( f(u) = u^5 \), og den indre funktion er \( g(x) = x^3 + 2x \). Deres afledte er \( f'(u) = 5u^4 \) og \( g'(x) = 3x^2 + 2 \).
Vi anvender kædereglen:
\[
\begin{aligned}
h'(x) &= f'(g(x)) \cdot g'(x) \\
&= 5(x^3 + 2x)^4 \cdot (3x^2 + 2)
\end{aligned}
\]
\end{solution}

\begin{example}

    Differentiér \(\left(3 x^2-5\right)^3\).
\end{example}

\begin{solution}
    Første skridt er altid at se, at udtrykket er en sammensat funktion, og derefter at opdele det i dets ydre og indre del. I dette eksempel er den ydre funktion $(\,\cdot\,)^{3}$ med den afledte $3(\,\cdot\,)^{2}$, og den indre funktion er $3x^{2}-5$ med den afledte $6x$. Reglen for sammensatte funktioner giver derfor

    \[
    \frac{d}{dx}\big(3x^{2}-5\big)^{3}
    = 3\big(3x^{2}-5\big)^{2} \times 6x
    = 18x\big(3x^{2}-5\big)^{2}.
    \]
    
    Alternativt kan vi indføre substitutionen \(u = 3x^{2}-5\) og skrive \(y = u^{3}\). Så er
    
    \[
    \frac{dy}{dx}
    = \frac{dy}{du} \times \frac{du}{dx}
    = 3u^{2} \times 6x
    = 18x\big(3x^{2}-5\big)^{2}.
    \]
    
\end{solution}

\begin{example}

    Find \(\frac{d y}{d x}\), hvis \(y=\sqrt{x^2+1}\).

\end{example}
\begin{solution}
    Den ydre funktion er \(\sqrt{\cdot}\) med den afledte \(\frac{1}{2\sqrt{\cdot}}\), og den indre funktion er \(x^2+1\) med den afledte \(2x\). Kædereglen giver derfor

    \[
    \frac{d y}{d x} = \frac{1}{2\sqrt{x^2+1}} \times 2x = \frac{x}{\sqrt{x^2+1}}.
    \]
\end{solution}

Kædereglen er den centrale mekanisme bag \textbf{backpropagation}-algoritmen, som bruges til at træne neurale netværk. En fejl i et netværks output er en sammensætning af mange indlejrede funktioner (lagene). Backpropagation bruger kædereglen gentagne gange til at beregne den afledte af fejlen med hensyn til hver vægt i netværket og fortæller os dermed, hvordan hver vægt skal justeres for at forbedre modellen.

\begin{example} En opgave med flere regler

Find den afledte af \( h(x) = \frac{x^2}{(3x+1)^4} \).
\end{example}

\begin{solution}
Opgaven kræver kvotientreglen, og nævneren kræver kædereglen.
Lad \( f(x) = x^2 \) og \( g(x) = (3x+1)^4 \).
Så er \( f'(x) = 2x \). For at finde \( g'(x) \) bruger vi kædereglen: \( g'(x) = 4(3x+1)^3 \cdot 3 = 12(3x+1)^3 \).
Nu anvender vi kvotientreglen:
\[
\begin{aligned}
h'(x) &= \frac{f'(x)g(x) - f(x)g'(x)}{[g(x)]^2} \\
&= \frac{(2x)(3x+1)^4 - (x^2)(12(3x+1)^3)}{((3x+1)^4)^2} \\
&= \frac{(2x)(3x+1) - 12x^2}{(3x+1)^5} \quad \text{(efter at have sat \((3x+1)^3\) uden for parentes og forkortet)} \\
&= \frac{6x^2 + 2x - 12x^2}{(3x+1)^5} = \frac{2x - 6x^2}{(3x+1)^5}
\end{aligned}
\]
\end{solution}

\subsection*{Oversigt over differentiationsreglerne}
\label{sec:ch10-common}

Differentiationsreglerne er samlet i følgende sætning.

\begin{theorem}{Differentiationsregler}

    \begin{enumerate}
        \item \textbf{Konstantreglen:} Hvis $c$ er en konstant, så er $\dfrac{d}{dx}c=0$. \vspace{0.5em}
        \item \textbf{Potensreglen:} For ethvert reelt $n$ gælder $\dfrac{d}{dx}x^n = n x^{n-1}$. \vspace{0.5em}
        \[
        \frac{d}{dx}x^n = n x^{n-1}.
        \]
        \item \textbf{Skalarreglen:} Hvis $c$ er en konstant, så er $\dfrac{d}{dx}[c \cdot f(x)] = c \cdot f'(x)$. \vspace{0.5em}
        \item \textbf{Linearitet:} For differentiable $f,g$ og konstanter $a,b$ gælder $\dfrac{d}{dx}(a f + b g) = a f' + b g'$. \vspace{0.5em}
        \item \textbf{Produktreglen:} Hvis $f,g$ er differentiable, så er $\dfrac{d}{dx}(f g) = f'g + fg'$. \vspace{0.5em}
        \item \textbf{Kvotientreglen:} Hvis $f,g$ er differentiable, og $g(x)\neq0$, så er $\dfrac{d}{dx}\left(\frac{f}{g}\right) = \frac{f'g - fg'}{g^2}$. \vspace{0.5em}
        \item \textbf{Kædereglen:} Hvis $f,g$ er differentiable, så er $\dfrac{d}{dx}(f(g(x))) = f'(g(x))g'(x)$. \vspace{0.5em}
    \end{enumerate}

\end{theorem}

Reglerne dækker ikke alle funktioner. Den afledte af $e^{x}$ er for eksempel $e^{x}$, mens den afledte af $\ln x$ er $\dfrac{1}{x}$, og ingen af dem passer ind i de mønstre, vi har beskrevet indtil nu. De mest almindelige differentiable funktioner, der falder uden for de tidligere regler, er anført nedenfor.


\begin{itemize}
    \item $\dfrac{d}{dx}e^x = e^x$\vspace{0.5em}
    \item $\dfrac{d}{dx}\ln x = \dfrac{1}{x}$\vspace{0.5em}
    \item $\dfrac{d}{dx}\sin x = \cos x$\vspace{0.5em}
    \item $\dfrac{d}{dx}\cos x = -\sin x$
\end{itemize}

\begin{example}

    Differentiér \(\ln \left(x^2+3 x+1\right)\).
\end{example}

\begin{solution} Vi løser opgaven ved hjælp af kædereglen og vores viden om den afledte af \(\ln x\).

    \[
    \begin{aligned}
    \frac{d}{d x} \ln\left(x^2+3 x+1\right) & =\frac{d}{d x}\left(\ln u\right) \quad\left(\text { hvor } u=x^2+3 x+1\right) \\
    & =\frac{d}{d u}\left(\ln u\right) \times \frac{d u}{d x} \quad(\text { ifølge kædereglen }) \\
    & =\frac{1}{u} \times \frac{d u}{d x} \\
    & =\frac{1}{x^2+3 x+1} \times \frac{d}{d x}\left(x^2+3 x+1\right) \\
    & =\frac{1}{x^2+3 x+1} \times(2 x+3) \\
    & =\frac{2 x+3}{x^2+3 x+1}
    \end{aligned}
    \]

\end{solution}


\begin{example}
    
Find \(\frac{d}{d x}\left(e^{3 x^2}\right)\).
\end{example}
\begin{solution} Det er en anvendelse af kædereglen sammen med vores viden om den afledte af \(e^x\).

\[
\begin{aligned}
\frac{d}{d x}\left(e^{3 x^2}\right) & =\frac{d e^u}{d x} \quad \text { hvor } u=3 x^2 \\
& =\frac{d e^u}{d u} \times \frac{d u}{d x} \quad \text { ifølge kædereglen } \\
& =e^u \times \frac{d u}{d x} \\
& =e^{3 x^2} \times \frac{d}{d x}\left(3 x^2\right) \\
& =6 x e^{3 x^2}
\end{aligned}
\]
\end{solution}


\begin{example}

Find \(\frac{d}{d x}\left(e^{x^2+2 x}\right)\).
\end{example}
\begin{solution} Igen bruger vi vores viden om den afledte af \(e^x\) sammen med kædereglen.

\[
\begin{aligned}
\frac{d}{d x}\left(e^{x^2+2 x}\right) & =\frac{d e^{u}}{d x} \quad\left(\text { hvor } u=x^2+2 x\right) \\
& =e^{u} \times \frac{d u}{d x} \quad \text { (ifølge kædereglen) } \\
& =e^{x^2+2 x} \times \frac{d}{d x}\left(x^2+2 x\right) \\
& =\left(2x+2\right)e^{x^2+2 x}
\end{aligned}
\]
\end{solution}


\begin{example}

Differentiér \(\ln \left(2 x^3+5 x^2-3\right)\).
\end{example}
\begin{solution}
Vi løser opgaven ved hjælp af kædereglen og vores viden om den afledte af \(\ln x\).

\[
\begin{aligned}
\frac{d}{d x} \ln \left(2 x^3+5 x^2-3\right) & =\frac{d \ln u}{d x} \quad \text { (hvor } u=\left(2 x^3+5 x^2-3\right) \\
& =\frac{d \ln u}{d u} \times \frac{d u}{d x} \quad \text { (ifølge kædereglen) } \\
& =\frac{1}{u} \times \frac{d u}{d x} \\
& =\frac{1}{2 x^3+5 x^2-3} \times \frac{d}{d x}\left(2 x^3+5 x^2-3\right) \\
& =\frac{1}{2 x^3+5 x^2-3} \times\left(6 x^2+10 x\right) \\
& =\frac{6 x^2+10 x}{2 x^3+5 x^2-3}
\end{aligned}
\]
\end{solution}

Der findes to genveje til at differentiere funktioner med eksponentialfunktioner og logaritmer. De fire eksempler ovenfor gav

\[
\begin{aligned}
\frac{d}{d x}\left(\ln\left(x^2+3 x+1\right)\right) & =\frac{2 x+3}{x^2+3 x+1} \\
\frac{d}{d x}\left(e^{3 x^2}\right) & =6 x e^{3 x^2} \\
\frac{d}{d x}\left(e^{x^2+2 x}\right) & =\left(2x+2\right)e^{x^2+2 x} \\
\frac{d}{d x}\left(\ln\left(2 x^3+5 x^2-3\right)\right) & =\frac{6 x^2+10 x}{2 x^3+5 x^2-3}
\end{aligned}
\]


Eksemplerne antyder de generelle regler

\[
\begin{aligned}
\frac{d}{d x}\left(e^{f(x)}\right) & =f^{\prime}(x) e^{f(x)} \\
\frac{d}{d x}(\ln f(x)) & =\frac{f^{\prime}(x)}{f(x)}
\end{aligned}
\]


Reglerne følger af kædereglen og af, at \(\frac{d}{dx}e^x=e^x\) og \(\frac{d}{dx}\ln x=\frac{1}{x}\) for $x>0$. De kan gøre differentiationen hurtigere, men hvis du glemmer dem, kan de genskabes ved hjælp af kædereglen som i eksemplerne ovenfor.


\section{Anvendelser af afledte}
\label{sec:ch10-applications}

Matematikkens udvikling hører til blandt den menneskelige tænknings største bedrifter, og fremkomsten af infinitesimalregningen -- både differential- og integralregning -- markerer et vendepunkt i den historie. Især differentialregning har utallige praktiske anvendelser inden for naturvidenskab og ingeniørfag, alt for mange til at opregne. Dens betydning afspejles i, at stort set alle kvantitative discipliner bygger på dens metoder.

I elementær matematik bruges differentialregning primært på to områder: til at analysere og skitsere kurver og til at løse optimeringsproblemer. I dette afsnit giver vi en kort introduktion til, hvordan differentialregning anvendes til optimering.


\subsection*{Stationære punkter -- idéen bag optimering}
Lad os vende tilbage til vores metafor med bjergkæden. Forestil dig, at du står et sted mellem tinder og dale, måske endda med lukkede øjne eller indhyllet i tåge. Mens du vandrer rundt, bemærker du, at når du går op ad bakke eller ned ad bakke, ved du, at du ikke er på den højeste top eller i den laveste dal -- du er enten stadig på vej op eller stadig på vej ned.

Men på toppen af en tinde er der et kort, plant stykke, hvor hældningen under dine fødder er nul. I bunden af en dal er der også et plant sted. Hvis du befinder dig på helt flad grund -- i det mindste et øjeblik -- har du et fingerpeg: Du er måske på en top, i en dal eller måske ved et blidt vendepunkt langs en bjergryg. I en bjergkædes store terræn er det kun på disse plane steder, hvor hældningen forsvinder, at du overhovedet kan være i et ekstremumspunkt: et maksimum eller et minimum.

Det motiverer et vigtigt skridt i optimering: Indre punkter, hvor en differentiabel funktion har et lokalt maksimum eller minimum, skal have afledt lig med nul. Disse steder kaldes stationære punkter. Et fuldstændigt optimeringsproblem kan også kræve, at man undersøger endepunkter og punkter, hvor den afledte ikke eksisterer.

\begin{definition}
     For en funktion $y=f(x)$ kaldes de punkter på grafen, hvor grafen har hældningen nul, for stationære punkter. Stationære punkter er med andre ord der, hvor $f^{\prime}(x)=0$.
    \end{definition}

For at finde en funktions stationære punkter differentierer vi, sætter den afledte lig med nul og løser ligningen.

\begin{example}

    Find de stationære punkter for funktionen $f(x)=2 x^{3}+3 x^{2}-12 x+17$.
    \label{ex:stationary-points-1}
\end{example}


\begin{solution}
    
    $f^{\prime}(x)=6 x^{2}+6 x-12$. Vi sætter $f^{\prime}(x)=0$ og løser:


\[
\begin{aligned}
6 x^{2}+6 x-12 & =0 \\
x^{2}+x-2 & =0 \\
(x-1)(x+2) & =0 \\
x & =1,-2
\end{aligned}
\]

Det giver de værdier af $x$, hvor funktionen $f$ er stationær. De tilhørende funktionsværdier findes ved at indsætte 1 og -2 i funktionen.

De er $f(1)=2 \times 1^{3}+3 \times 1^{2}-12 \times 1+17=10$ og 
$f(-2)=2 \times(-2)^{3}+3 \times(-2)^{2}-12 \times(-2)+17=37$. De stationære punkter er derfor $(1,10)$ og $(-2,37)$.
\end{solution}


\begin{example}

    Find de stationære punkter for funktionen $g(t)=e^{t^{2}}$.
\end{example}

\begin{solution}
    Ved at differentiere og sætte den afledte lig med nul får vi ligningen $g^{\prime}(t)=2 t e^{t^{2}}=0$. Da $e^{t^{2}}$ aldrig er nul, er den eneste løsning til ligningen der, hvor $2 t=0$, dvs. $t=0$. Ved indsættelse i forskriften for $g$ får vi funktionsværdien $g(0)=e^{0^{2}}=1$. Det stationære punkt er altså $(0,1)$.
\end{solution}

\subsubsection*{Typer af stationære punkter}
Hvis vi vender tilbage til metaforen med bjergkæden, kan stationære punkter opfattes som de særlige steder, hvor jorden et øjeblik føles helt plan under dine fødder. Nogle gange sker det helt oppe på toppen af en bakke -- hvor du, uanset hvilken retning du rækker ud i, kun mærker skråninger nedad foran dig. Det kaldes et lokalt maksimum: toppen i dine umiddelbare omgivelser, men ikke nødvendigvis den højeste tinde i hele bjergkæden. Andre gange ligger det plane sted i bunden af en dal -- alt omkring dig skråner opad. Det kaldes et lokalt minimum: det laveste sted i dit lokale nærområde, selv om der findes dybere dale andre steder. Betegnelsen »lokal« understreger, at vi taler om det højeste eller laveste punkt i nærheden, ikke globalt. Du vil måske også høre betegnelserne »relativt maksimum« og »relativt minimum« brugt om dem. \autoref{fig:local-extrema} viser en funktion som et landskab med både en bakketop (lokalt maksimum) og en dalbund (lokalt minimum). Bemærk, at hældningen er nul i hvert af disse punkter -- præcis som på et plant stykke i bjergene.

% Figure: Local maximum and local minimum (example with f(x)=x^3 - 3x)
% --- Ensure these packages are in your preamble ---
% \usepackage{pgfplots}
% \usepackage{xcolor}
% \pgfplotsset{compat=1.17} % Or your current version

% --- Ensure your custom color is defined ---
% \definecolor{mseViaBlue}{HTML}{1933CC}

\begin{figure}[htbp]
    \centering
    \begin{tikzpicture}
        \begin{axis}[
            axis lines=middle,
            xlabel=$x$,
            ylabel=$y$,
            width=12cm, height=8cm, % Slightly adjusted size
            xmin=-2.5, xmax=2.5,
            ymin=-3.5, ymax=3.5,
            xtick={-2, -1, 0, 1, 2},
            ytick={-3, -2, -1, 0, 1, 2, 3},
            grid=major,
            grid style={dashed, gray!30},
            legend pos=outer north east,
            clip=false % Allows labels to go outside the axis box
        ]
        % Plot the curve f(x) = x^3 - 3x, specifying the domain
        \addplot[
            domain=-2.2:2.2, % Explicitly set the domain
            samples=100,     % Ensure a smooth curve
            smooth,
            very thick,
            mseViaBlue
        ] {x^3 - 3*x};
        \addlegendentry{\(f(x)=x^3-3x\)}

        % Mark the stationary points: (-1, 2) and (1, -2)
        \addplot[
            only marks,
            mark=*,              % A solid circle is often clearer
            mark size=2.5pt,
            color=red            % Make the points stand out
        ] coordinates {(-1,2) (1,-2)};
        \addlegendentry{Stationære punkter}

        % Add labels for the points using nodes
        \node[above right, color=red, font=\small] at (axis cs:-1, 2) {Lokalt maksimum};
        \node[below right, color=red, font=\small] at (axis cs:1, -2) {Lokalt minimum};
        \end{axis}
    \end{tikzpicture}
    \caption{En kurve med et lokalt maksimum i \(x=-1\) og et lokalt minimum i \(x=1\).}
    \label{fig:local-extrema}
\end{figure}

Lokale maksima og lokale minima er ikke de eneste typer af stationære punkter. Der findes en tredje slags. \autoref{fig:stationary-inflection} viser et stationært punkt, der hverken er et lokalt maksimum eller et lokalt minimum. Denne type stationært punkt kaldes et stationært vendepunkt eller blot et \textit{vendepunkt}.

% --- Ensure these packages are in your preamble ---
% \usepackage{pgfplots}
% \usepackage{xcolor}
% \pgfplotsset{compat=1.17} % Or your current version

% --- Ensure your custom color is defined ---
% \definecolor{mseViaBlue}{HTML}{1933CC}

\begin{figure}[htbp]
    \centering
    \begin{tikzpicture}
        \begin{axis}[
            axis lines=middle,
            xlabel=$x$,
            ylabel=$y$,
            width=12cm, height=8cm,
            xmin=-1.5, xmax=3.5,
            ymin=-4.5, ymax=4.5,
            xtick={-1, 0, 1, 2, 3},
            ytick={-4, -2, 0, 2, 4},
            grid=major,
            grid style={dashed, gray!30},
            legend pos=outer north east,
            legend cell align={left},
            clip=false
        ]
        % Plot the curve f(x) = (x-1)^3
        \addplot[
            domain=-0.8:2.8, % Explicitly set the domain
            samples=100,      % Ensure a smooth curve
            smooth,
            very thick,
            mseViaBlue
        ] {(x-1)^3};
        \addlegendentry{\(f(x)=(x-1)^3\)}

        % Mark the stationary point of inflection at x=1
        \addplot[
            only marks,
            mark=*,              % Use a solid circle
            mark size=2.5pt,
            color=red
        ] coordinates {(1,0)};
        \addlegendentry{Stationært punkt}

        % Add a label for the point using a node and an arrow
        \node[color=red, font=\small] (label_node) at (axis cs:2.5, 2) {Stationært vendepunkt};
        \draw[->, color=red, thick] (label_node) -- (axis cs:1.1, 0.1);

        % Optional: Add a tangent line at the inflection point to show it's horizontal
        \addplot[
            domain=0:2,
            dashed,
            color=black,
            thin
        ] {0};
        \addlegendentry{Tangent i x=1}

        \end{axis}
    \end{tikzpicture}
    \caption{En funktion med et stationært vendepunkt i \(x=1\). Tangenten i punktet er vandret, men punktet er hverken et lokalt maksimum eller et lokalt minimum.}
    \label{fig:stationary-inflection}
\end{figure}

\subsection*{Første afledtes test}
Lad os vende tilbage til \autoref{ex:stationary-points-1}. For funktionen $f(x)=2 x^{3}+3 x^{2}-12 x+17$ fandt vi stationære punkter i $(1,10)$ og $(-2,37)$. Det naturlige spørgsmål er: Hvilken type stationære punkter er det? Uden en skitse af grafen er det ikke umiddelbart klart, om de svarer til lokale maksima, lokale minima eller stationære vendepunkter. En nøjagtig graf ville besvare spørgsmålet, men den kan kræve en betydelig indsats. I stedet søger vi en metode, der gør det muligt at klassificere et stationært punkt uden at tegne hele funktionen.

Der findes flere fremgangsmåder, men i dette afsnit fokuserer vi på én teknik, kendt som \emph{første afledtes test}. Idéen er at undersøge funktionens opførsel umiddelbart til venstre og umiddelbart til højre for det stationære punkt.

For at forstå princippet kan du forestille dig en person, der står et sted i et landskab, hvor jorden er plan. Uden at kunne se omgivelserne vil personen afgøre, om stedet er toppen af en bakke, bunden af en dal eller ingen af delene. En måde at afgøre det på er at tage et lille skridt tilbage og mærke hældningen og derefter tage et lille skridt frem og mærke hældningen igen.

Hvis jorden skråner opad bagude og nedad forude, er det plane punkt toppen af en bakke, svarende til et lokalt maksimum. Hvis jorden skråner nedad bagude og opad forude, er punktet bunden af en dal, svarende til et lokalt minimum. Hvis hældningen har samme fortegn på begge sider, er punktet ikke et lokalt ekstremum; der kræves yderligere analyse for at afgøre, om det er et vendepunkt.

Denne tankegang indfanger essensen af første afledtes test: Ved at analysere fortegnet for \(f^{\prime}(x)\) umiddelbart på hver side af et stationært punkt kan vi bestemme punktets karakter uden at være afhængige af en graf.

% --- Ensure these packages are in your preamble ---
% \usepackage{pgfplots}
% \usepackage{xcolor}
% \pgfplotsset{compat=1.17}

% --- Ensure your custom color is defined ---
% \definecolor{mseViaBlue}{HTML}{1933CC}

\begin{figure}[htbp]
    \centering
    \begin{tikzpicture}
        \begin{axis}[
            axis lines=middle,
            xlabel=$x$,
            ylabel=$y$,
            width=12cm, height=8cm,
            xmin=-2.5, xmax=2.5,
            ymin=-3.5, ymax=3.5,
            xtick={-2, -1, 0, 1, 2},
            ytick={-3, -2, -1, 0, 1, 2, 3},
            grid=major,
            grid style={dashed, gray!30},
            legend pos=outer north east,
            clip=false
        ]
        % Plot the curve f(x) = x^3 - 3x
        \addplot[
            domain=-2.2:2.2, samples=100,
            smooth, very thick, mseViaBlue
        ] {x^3 - 3*x};
        \addlegendentry{\(f(x)=x^3-3x\)}

        % -- Local Maximum at x = -1 --
        \coordinate (max_point) at (axis cs:-1, 2);
        % Tangent with positive slope (f' > 0)
        \draw[green!60!black, thick] (axis cs:-1.8, 0.5) -- (axis cs:-1.2, 3.5);
        \node[above, color=green!60!black] at (axis cs:-1.7, 2.2) {\(f' > 0\)};
        % Tangent with negative slope (f' < 0)
        \draw[red, thick] (axis cs:-0.8, 3.5) -- (axis cs:-0.2, 0.5);
        \node[above, color=red] at (axis cs:-0.3, 2.2) {\(f' < 0\)};
        % Horizontal tangent (f' = 0)
        \draw[black, dashed] (axis cs:-1.5, 2) -- (axis cs:-0.5, 2);
        \node[above] at (max_point) {\(f' = 0\)};

        % -- Local Minimum at x = 1 --
        \coordinate (min_point) at (axis cs:1, -2);
        % Tangent with negative slope (f' < 0)
        \draw[red, thick] (axis cs:0.2, -0.5) -- (axis cs:0.8, -3.5);
        \node[below, color=red] at (axis cs:0.3, -2.5) {\(f' < 0\)};
        % Tangent with positive slope (f' > 0), mirrored around x = 1 to match the negative slope tangent
        \draw[green!60!black, thick] (axis cs:1.2, -3.5) -- (axis cs:1.8, -0.5);
        \node[below, color=green!60!black] at (axis cs:1.7, -2.5) {\(f' > 0\)};
        % Horizontal tangent (f' = 0)
        \draw[black, dashed] (axis cs:0.5, -2) -- (axis cs:1.5, -2);
        \node[below] at (min_point) {\(f' = 0\)};

        % Mark the stationary points
        \fill[red] (max_point) circle (2.5pt);
        \fill[red] (min_point) circle (2.5pt);

        \end{axis}
    \end{tikzpicture}
    \caption{Illustration af første afledtes test. I et lokalt maksimum skifter den afledte fra positiv til negativ. I et lokalt minimum skifter den fra negativ til positiv.}
    \label{fig:first-derivative-test}
\end{figure}

% --- Ensure these packages are in your preamble ---
% \usepackage{pgfplots}
% \usepackage{xcolor}
% \pgfplotsset{compat=1.17}

% --- Ensure your custom color is defined ---
% \definecolor{mseViaBlue}{HTML}{1933CC}

\begin{figure}[htbp]
    \centering
    \begin{tikzpicture}
        \begin{axis}[
            axis lines=middle,
            xlabel=$x$,
            ylabel=$y$,
            width=12cm, height=8cm,
            xmin=-2.5, xmax=2.5,
            ymin=-2.5, ymax=2.5,
            xtick={-2, -1, 0, 1, 2},
            ytick={-2, -1, 0, 1, 2},
            grid=major,
            grid style={dashed, gray!30},
            legend pos=outer north east,
            clip=false
        ]
        % Plot the curve f(x) = x^3 / 2
        \addplot[
            domain=-2:2, samples=100,
            smooth, very thick, mseViaBlue
        ] {x^3 / 2};
        \addlegendentry{\(f(x)=\frac{1}{2}x^3\)}

        % -- Stationary Point of Inflection at x = 0 --
        \coordinate (inflection_point) at (axis cs:0, 0);

        % Tangent with positive slope to the left (f' > 0)
        \draw[green!60!black, thick] (axis cs:-2, -2) -- (axis cs:-0.5, 0);
        \node[above, color=green!60!black] at (axis cs:-1.5, -1) {\(f' > 0\)};

        % Tangent with positive slope to the right (f' > 0)
        \draw[green!60!black, thick] (axis cs:0.5, 0) -- (axis cs:2, 2);
        \node[below, color=green!60!black] at (axis cs:1.5, 1) {\(f' > 0\)};

        % Horizontal tangent at the inflection point (f' = 0)
        \draw[black, dashed] (axis cs:-1, 0) -- (axis cs:1, 0);
        \node[above, anchor=south west] at (inflection_point) {\(f' = 0\)};

        % Mark the stationary point
        \fill[red] (inflection_point) circle (2.5pt);

        \end{axis}
    \end{tikzpicture}
    \caption{Illustration af et stationært vendepunkt. Den afledte er nul i \(x=0\), men positiv på begge sider, så punktet er hverken et lokalt maksimum eller et lokalt minimum.}
    \label{fig:stationary-inflection-test}
\end{figure}

Vi kan sammenfatte disse overvejelser i følgende sætning:

\begin{theorem}[Første afledtes test]
    Lad \(c\) være et kritisk punkt for en kontinuert funktion \(f\).
    \begin{itemize}
        \item Hvis \(f'(x)\) skifter fra negativ til positiv i \(c\), har \(f\) et \textbf{lokalt minimum} i \(c\).
        \item Hvis \(f'(x)\) skifter fra positiv til negativ i \(c\), har \(f\) et \textbf{lokalt maksimum} i \(c\).
        \item Hvis \(f'(x)\) ikke skifter fortegn i \(c\), har \(f\) intet lokalt ekstremum i \(c\). Første afledtes test kan ikke alene afgøre, om $c$ er et vendepunkt.
    \end{itemize}

    \end{theorem}
Testen er illustreret i \autoref{fig:first-derivative-test} og \autoref{fig:stationary-inflection-test}.


\subsection*{Optimering}
Nu hvor vi har lagt fundamentet, kan vi gå i gang med nogle optimeringsproblemer. Når vi vil maksimere en funktion $f(x)$ på et bestemt interval for $x$, er målet at finde den største værdi, som $f(x)$ antager på intervallet. Det er vigtigt at være opmærksom på, at denne maksimumsværdi ikke altid ligger i et stationært punkt. Se på \autoref{fig:restricted-domain-extrema}, hvor vi søger funktionens største og mindste værdi på intervallet $2 \leq x \leq 7$. På intervallet har funktionen to stationære punkter -- det ene svarer til et lokalt maksimum, det andet til et lokalt minimum. Funktionens største værdi på intervallet antages dog faktisk i endepunktet $x=7$, som ikke er et stationært punkt. Funktionen antager denne værdi, simpelthen fordi større $x$-værdier uden for intervallet ikke tages i betragtning. Mindsteværdien ligger derimod i dette eksempel i et stationært punkt inden for intervallet. Med denne forståelse kan vi nu forklare de systematiske trin til at finde en funktions maksimum eller minimum i et givet område.

% --- Ensure these packages are in your preamble ---
% \usepackage{pgfplots}
% \usepackage{xcolor}
% \pgfplotsset{compat=1.17}
% \usetikzlibrary{patterns} % For the shaded region

% --- Ensure your custom color is defined ---
% \definecolor{mseViaBlue}{HTML}{1933CC}

\begin{figure}[htbp]
    \centering
    \begin{tikzpicture}
        \begin{axis}[
            axis lines=middle, % <-- Change from 'left' to 'middle'
            xlabel=$x$,
            ylabel=$y$,
            width=10cm, height=7cm,
            xmin=0, xmax=8.5,
            ymin=-5, ymax=75,
            xtick={0, 1, 2, 3, 4, 5, 6, 7, 8},
            ytick={0, 25, 50, 75},
            grid=major,
            grid style={dashed, gray!30},
            legend style={
                at={(0.5, 0.97)}, % Position near the top center
                anchor=north,     % Anchor the top-center of the legend box
                legend cell align={left}
            },
            clip=false,
            % Add setting to move x axis to y=0
            axis x line=bottom, % Ensures x axis is shown
            axis y line=middle, % y axis through x=0
            % Customize placement:
            x axis line style={-}, % Solid x axis
            y axis line style={-}, % Solid y axis
            enlarge y limits={abs=0.5}, % Provide space above/below
            % Custom ticks at y=0 (x axis at y=0)
            at={(0,0)}
        ]
        % Define the function path for shading
        \addplot[
            name path=curve,
            domain=1.5 :7.5, samples=100,
            smooth, very thick, mseViaBlue
        ] {2*(x^3 - 12*x^2 + 45*x) - 99};
        \addlegendentry{\(f(x)\)}

        % -- Highlight the restricted domain [2, 7] --
        % Create paths for the vertical boundaries
        \path[name path=left_bound] (axis cs:2, -5) -- (axis cs:2, 45);
        \path[name path=right_bound] (axis cs:7, -5) -- (axis cs:7, 45);
        % Shade the area between x=2 and x=7
        \addplot[gray!20] fill between[of=left_bound and right_bound];
        \addlegendentry{Interval [2, 7]}

        % -- Mark and Label the Points of Interest --
        % Local Maximum (at x=3)
        \coordinate (local_max) at (axis cs:3, 9);
        \fill[blue] (local_max) circle (2.5pt);
        \node[above, blue, font=\small] at (local_max) {Lokalt maksimum};

        % Absolute Minimum (at x=5, which is also a local min)
        \coordinate (abs_min) at (axis cs:5, 1);
        \fill[red] (abs_min) circle (2.5pt);
        \node[below, red, font=\small] at (abs_min) {Globalt minimum};

        % Absolute Maximum (at endpoint x=7)
        \coordinate (abs_max) at (axis cs:7, 41);
        \fill[red] (abs_max) circle (2.5pt);
        \node[above, red, font=\small] at (abs_max) {Globalt maksimum};

        % Left Endpoint (at x=2)
        \coordinate (left_end) at (axis cs:2, 1);
        \fill[black] (left_end) circle (2.5pt);

        \end{axis}
    \end{tikzpicture}
    \caption{På det begrænsede interval \([2, 7]\) ligger det globale maksimum i et endepunkt (\(x=7\)) og ikke i det lokale maksimum. Det globale minimum ligger i et indre stationært punkt (\(x=5\)).}
    \label{fig:restricted-domain-extrema}
\end{figure}

\subsection*{Placeringen af maksima og minima}
En funktion $f(x)$ har ikke nødvendigvis en største eller mindste værdi i et bestemt område af $x$-værdier. Men hvis de findes, skal maksimums- og minimumsværdien ligge et af tre steder:

\begin{enumerate}
  \item I endepunkterne (hvis de findes) af det område, vi betragter.
  \item Inde i området i et stationært punkt.
  \item Inde i området i et punkt, hvor den afledte ikke eksisterer.
\end{enumerate}

\begin{remark}
    

\begin{enumerate}
  \item Det er let at finde tilfælde, hvor en funktion ikke antager et maksimum eller minimum inden for et bestemt område. Funktionen $f(x) = x$ har for eksempel hverken maksimum eller minimum på intervallet $-\infty < x < \infty$; grafen vokser ubegrænset, når $x$ vokser. Med hensyn til punkt 1 ovenfor skal du bemærke, at området $-\infty < x < \infty$ faktisk ikke har nogen endepunkter. Et andet eksempel er området $x \geq 1$, som kun har ét endepunkt, nemlig $x = 1$.
  \item Med hensyn til punkt 3 behandler denne tekst ikke situationer, hvor den afledte ikke eksisterer i visse punkter. Husk dog, at sådanne punkter kan forekomme, og at et maksimum eller minimum kan ligge netop dér. Flere detaljer findes i litteratur om videregående infinitesimalregning.
\end{enumerate}
\end{remark}

Nu hvor vi ved præcis, hvor maksima og minima kan ligge, kan vi give en fremgangsmåde til at finde dem.

\subsubsection*{Fremgangsmåde til at finde en funktions største eller mindste værdi.}
\begin{enumerate}
  \item Find endepunkterne for det område, vi betragter (hvis der er nogen).
  \item Find alle stationære punkter i området.
  \item Find alle punkter i området, hvor den afledte ikke eksisterer.
  \item Indsæt hvert af disse punkter i funktionen, og se, hvilket der giver den største (eller mindste) funktionsværdi.
\end{enumerate}

\begin{example}

    Find den mindste og den største værdi af funktionen $f(x)=x^{2} e^{x}$ for $-4 \leq x \leq 1$.
\end{example}

\begin{solution} Vi følger fremgangsmåden ovenfor. Endepunkterne er -4 og 1. Ved differentiation får vi $f^{\prime}(x)=x^{2} e^{x}+2 x e^{x}=x(x+2) e^{x}$. Ved at sætte $f^{\prime}(x)=0$ og løse ligningen får vi stationære punkter i $x=0$ og $x=-2$. Der er ingen punkter, hvor den afledte ikke eksisterer. Den største og den mindste værdi findes derfor i et af punkterne $x=-4,-2,0,1$. Ved indsættelse får vi $f(-4) \approx 0.29, f(-2) \approx 0.54, f(0)=0$ og $f(1)=e \approx 2.7$. Den største værdi antages derfor i $x=1$ og er lig med $e$, og den mindste værdi antages i $x=0$ og er 0.
\end{solution}

\begin{example}

    Find den største og den mindste værdi af funktionen $g(t)=\frac{1}{3} t^{3}-t+2$ for $0 \leq t \leq 3$.
\end{example}

\begin{solution} Endepunkterne er $t=0$ og $t=3$. Ved at differentiere og sætte lig med nul får vi $g^{\prime}(t)=t^{2}-1=(t-1)(t+1)=0$, så de stationære punkter ligger i $t=-1,1$. Da -1 ikke ligger i området, er de mulige placeringer af maksimum og minimum $t=0,1,3$. Ved indsættelse i $g$ får vi $g(0)=2, g(1)=\frac{4}{3}$ og $g(3)=8$. Den største værdi er derfor $g(3)=8$, og den mindste værdi er $g(1)=\frac{4}{3}$.
\end{solution}

\begin{example}

    Antag, at du tuner en maskinlæringsmodel og vil finde den læringsrate \( \alpha \) (mellem \(0\) og \(1\)), der minimerer tabsfunktionen \( L(\alpha) = (\alpha - 0.3)^2 + 1 \). Hvilken værdi af \( \alpha \) giver det mindste tab?
    \end{example}
    
    \begin{solution} Vi anvender fremgangsmåden til at finde globale ekstrema på det lukkede interval \( [0, 1] \). Den funktion, der skal minimeres, er tabsfunktionen \( L(\alpha) = (\alpha - 0.3)^2 + 1 \).
    
    \begin{enumerate}
        \item \textbf{Find endepunkterne:} Det område, vi betragter, er \( 0 \le \alpha \le 1 \), så endepunkterne er \( \alpha = 0 \) og \( \alpha = 1 \).
    
        \item \textbf{Find de stationære punkter:} Først finder vi den afledte af tabsfunktionen med hensyn til \( \alpha \):
        \[
        \frac{dL}{d\alpha} = \frac{d}{d\alpha} \left[ (\alpha - 0.3)^2 + 1 \right] = 2(\alpha - 0.3)
        \]
        Dernæst sætter vi den afledte lig med nul for at finde de stationære punkter:
        \[
        2(\alpha - 0.3) = 0 \implies \alpha = 0.3
        \]
        Det eneste stationære punkt er \( \alpha = 0.3 \), som ligger i vores interval \( [0, 1] \).
    
        \item \textbf{Find punkter, hvor den afledte ikke er defineret:} Den afledte \( 2(\alpha - 0.3) \) er en simpel lineær funktion og er defineret for alle værdier af \( \alpha \). Der er ingen sådanne punkter.
    
        \item \textbf{Sammenlign værdierne i kandidatpunkterne:} Vi udregner nu tabsfunktionen \( L(\alpha) \) i alle de kandidatpunkter, vi har fundet: endepunkterne og det stationære punkt.
        \begin{itemize}
            \item I det venstre endepunkt, \( \alpha = 0 \):
            \[ L(0) = (0 - 0.3)^2 + 1 = 0.09 + 1 = 1.09 \]
            \item I det stationære punkt, \( \alpha = 0.3 \):
            \[ L(0.3) = (0.3 - 0.3)^2 + 1 = 0 + 1 = 1.00 \]
            \item I det højre endepunkt, \( \alpha = 1 \):
            \[ L(1) = (1 - 0.3)^2 + 1 = (0.7)^2 + 1 = 0.49 + 1 = 1.49 \]
        \end{itemize}
    \end{enumerate}
    
    Ved at sammenligne værdierne ser vi, at det mindste tab er \( 1.00 \), som antages i det stationære punkt.
    
    \textbf{Konklusion:} Det mindste tab opnås, når læringsraten er \( \alpha = 0.3 \).
    \end{solution}

    \section{Den anden afledte: krumning og vendepunkter}

    Den første afledte, \(f'(x)\), fortæller os om en funktions hældning eller ændringshastighed. Den \textbf{anden afledte}, betegnet \(f''(x)\), er den afledte af den første afledte. Den måler, hvordan selve hældningen ændrer sig. Det giver en dybere indsigt i formen af en funktions graf.
    
    \begin{definition}[Den anden afledte]
    Den \textbf{anden afledte} af en funktion \( f \), betegnet \( f''(x) \), er funktionen
    \[
    f''(x) = \frac{d}{dx} \left( f'(x) \right)
    \]
    forudsat at grænseværdien eksisterer.
    \end{definition}
    
    Den vigtigste geometriske fortolkning af den anden afledte er \textbf{krumning}.
    
    \begin{itemize}
        \item Hvis \( f''(x) > 0 \) på et interval, er hældningen \(f'(x)\) voksende. Grafen bøjer opad som en kop, der holder på vand. Vi siger, at funktionen er \textbf{konveks} (krummer opad).
        \item Hvis \( f''(x) < 0 \) på et interval, er hældningen \(f'(x)\) aftagende. Grafen bøjer nedad som en kop, der spilder vand. Vi siger, at funktionen er \textbf{konkav} (krummer nedad).
    \end{itemize}
    
    Et punkt på grafen, hvor krumningen skifter, kaldes et \textbf{vendepunkt}. Det sker typisk der, hvor \(f''(x) = 0\).
    
    \begin{figure}[htbp]
        \centering
        \begin{tikzpicture}
            \begin{axis}[
                axis lines=middle,
                xlabel=$x$,
                ylabel=$y$,
                width=12cm, height=8cm,
                xmin=-2.5, xmax=2.5,
                ymin=-3, ymax=3,
                xtick={-2, -1, 0, 1, 2},
                ytick={-2, 0, 2},
                grid=major,
                grid style={dashed, gray!30},
                legend style={at={(0.03, 0.97)}, anchor=north west},
                clip=false
            ]
            % Plot the curve f(x) = x^3/2 - x
            \addplot[
                domain=-2.2:2.2, samples=100,
                smooth, very thick, mseViaBlue
            ] {x^3/2 - x};
            \addlegendentry{\(f(x)=\frac{1}{2}x^3 - x\)}
    
            % Mark the inflection point at x=0
            \coordinate (P) at (axis cs:0, 0);
            \fill[red] (P) circle (2.5pt);
            \node[above right, font=\small, color=red] at (P) {Vendepunkt i \(x=0\)};
    
            % Annotate Concavity
            \node[below, font=\small, color=orange] at (axis cs:-1.5, -1) {Konkav (\(f'' < 0\))};
            \node[above, font=\small, color=green!60!black] at (axis cs:1.5, 1) {Konveks (\(f'' > 0\))};
            
            \end{axis}
        \end{tikzpicture}
        \caption{En funktion med et område, hvor den er konkav (krummer nedad), efterfulgt af et område, hvor den er konveks (krummer opad). Overgangen sker i vendepunktet.}
        \label{fig:ch11-concavity}
    \end{figure}
    
    Den anden afledte giver os også et stærkt værktøj til at klassificere stationære punkter.
    
    \begin{theorem}[Anden afledtes test]
    Lad \(c\) være et stationært punkt for \(f\) (dvs. \(f'(c) = 0\)).
    \begin{itemize}
        \item Hvis \(f''(c) > 0\), har \(f\) et \textbf{lokalt minimum} i \(c\).
        \item Hvis \(f''(c) < 0\), har \(f\) et \textbf{lokalt maksimum} i \(c\).
        \item Hvis \(f''(c) = 0\), giver testen ingen konklusion.
    \end{itemize}
    \end{theorem}
